\begin{textsample}{POS Dim 2 – rj – Score 17.00 – rj/rj\_em\_3.txt}  \label{ex:f2_pos_216}
\textbf{LEGISLAÇÃO} \textbf{VIGENTE} E \textbf{CURRÍCULO} Seguindo os pressupostos da BNCC em relação a uma pedagogia das competências, cujo intuito é afirmar valores e estimular ações que contribuam para a transformação da sociedade, tornando -a mais humana, justa e, também, voltada para a preservação da natureza, o \textbf{Currículo} Referencial do Estado do Rio de Janeiro relaciona a \textbf{legislação} brasileira \textbf{vigente} às habilidades e competências a serem desenvolvidas no processo educativo.
De acordo com o Art. 9{\EmojiFont °} da Resolução SEEDUC Nº 6035, de 28 de \textbf{janeiro} de 2022, no âmbito de todo o \textbf{currículo} escolar, deverão ser inclusos princípios, aspectos e conteúdos relativos às Leis balizadoras abaixo listadas : - Lei Federal n{\EmojiFont °} 9.795/1999 : Política Nacional de Educação Ambiental ; - Resolução n{\EmojiFont °} 02/2012 : \textbf{Diretrizes} Curriculares Nacionais para a Educação Ambiental ; - Lei Federal n{\EmojiFont °} 12.068/2012 : Educação ambiental de forma integrada aos conteúdos obrigatórios, alterando a LDB 9394/1996 ; - Lei Federal n{\EmojiFont °} 10.639/2003 : \textbf{Diretrizes} e bases da educação nacional, para incluir no \textbf{currículo} oficial da Rede de Ensino a obrigatoriedade da temática"História e Cultura Afro-Brasileira"; - Lei Federal n{\EmojiFont °} 11.645/2008 : \textbf{Diretrizes} e bases da educação nacional, para incluir no \textbf{currículo} oficial da rede de ensino a obrigatoriedade da temática “ História e Cultura Afro-Brasileira e Indígena ” ; - Lei Federal 11.769/2008 : Altera a Lei no 9.394, de 20 de dezembro de 1996, Lei de \textbf{Diretrizes} e Bases da Educação, para dispor sobre a obrigatoriedade do ensino da música na educação básica ; - Lei Federal n{\EmojiFont °} 11.947/2009 : Dispõe sobre o atendimento da alimentação escolar e do Programa Dinheiro Direto na Escola aos alunos da educação básica ; altera as Leis nos 10.880, de 9 de junho de 2004, 11.273, de 6 de \textbf{fevereiro} de 2006, 11.507, de 20 de \textbf{julho} de 2007 ; revoga dispositivos da Medida Provisória no 2.178-36, de 24 de agosto de 2001, e a Lei no 8.913, de 12 de \textbf{julho} de 1994 ; e dá outras \textbf{providências} ; - Decreto Federal n{\EmojiFont °} 7.037/2009 : \textbf{Institui} o Programa Nacional dos Direitos Humano ; - Lei Federal n{\EmojiFont °} 13.010/2014 : Altera a Lei nº 8.069, de 13 de \textbf{julho} de 1990 ( Estatuto da Criança e do Adolescente ), para estabelecer o direito da criança e do \textbf{adolescente} de serem educados e cuidados sem o uso de castigos físicos ou de tratamento cruel ou degradante, e altera a Lei nº 9.394, de 20 de dezembro de 1996 ; - Lei Estadual n{\EmojiFont °} 3.721/2001 : Educação fiscal e financeira visando desenvolver uma consciência ética para consigo e para a comunidade ; - Lei Estadual 4.528/2005 : Estabelece as \textbf{Diretrizes} para a organização do Sistema de Ensino do Estado do Rio de Janeiro ; - Lei Federal n{\EmojiFont °} 10.741/2003 : Dispõe sobre o Estatuto do Idoso e dá outras \textbf{providências} e Lei Estadual n{\EmojiFont °} 6.820/2014 : Dispõe sobre a inclusão dos"Direitos dos Idosos"no \textbf{currículo} escolar das escolas do Estado do Rio de Janeiro na forma que menciona ; - Lei Federal n{\EmojiFont °} 13.005/2014 : Aprova o Plano Nacional de Educação - PNE e dá outras \textbf{providências}. - \textbf{Parecer} CEE/RJ n{\EmojiFont °} 60 de 14 de agosto de 2018 : Autoriza a Proposta Pedagógica, em caráter experimental por cinco anos, de atendimento semipresencial na Educação de Jovens e Adultos nas unidades escolares prisionais.
A construção coletiva e colaborativa do \textbf{Currículo} Referencial do Estado do Rio de Janeiro, com a participação dos professores da rede, procura alinhar os conceitos propostos pela Base Nacional Comum Curricular - BNCC com as especificidades encontradas na rede.
A proposta pedagógica do \textbf{currículo} permite uma dialogicidade na Educação de Jovens e Adultos, nas Unidades Escolares em Espaços de Privação de Liberdade presencial e semipresencial, nas Unidades Escolares em Espaços Socioeducativos e no Ensino Médio Regular.
Além disso, integra as características que marcam as diferenças regionais encontradas no amplo território do Estado do Rio de Janeiro que se refletem na vida social dos indivíduos.

% matched lemmas: adolescente, currículo, diretriz, fevereiro, instituir, janeiro, julho, legislação, parecer, providência, vigente
\end{textsample}
